% !TeX root = ../papers/nc-p3.tex
\section{異なる平面因子の一意性と横断する二軸}
\label{sec:distinct-factor-rigidity}
記号は第\ref{sec:square-apolar-ore}節を引き継ぐ。

\begin{theorem}[各因子の相手の一意性]\label{thm:unique-linear-factor-partner}
両対称の標準AS反転窓では全ての$0\ne x\in E^*$に対して
$\dim\{w:xw\in K_E\}\le1$であり、右側でも同じである。
従って反対階数8では両反対鉛筆が四次完全交叉を定める。
反対階数を問わず、添字保持同型の反対二位置での基底変換は射影的に一致する。
\end{theorem}
\begin{proof}
$X=\{w:xw\in K_E\}$とし、$f=\ker\beta(x,-)$を取る（$X\ne0$の場合）。
補題\ref{lem:reciprocal-apolar-factors}は$\mathbb P(X)$から
$\mathbb P(Y)$、$Y=\{v:fv\in K_F\}$への同型を与える。
$U=U_x$、$V=U^\perp$と置くと、全ての$0\ne w\in X,0\ne v\in Y$で
$U_w=V$、$V_v=U$である。$\dim X\ge2$なら$x$とは異なる因子を選べるので
$R=U\cap V$は高々一次元である。$\beta(X,Y)\subset R$で、
$w$の$Y$上の核は丁度その対応直線であるから$\dim X=\dim Y=2$、$\dim R=1$となる。
一方、各$0\ne v\in Y$で$V_v=U$なので
$E^*/X\to U/R$は同型である。これは二次元空間間の二行行列の線型鉛筆で、
行列式が射影直線上で消える点が存在し矛盾する。
階数8の独立な二二次式に共通因子があれば、その因子には二次元の相手がある。

次に、反対切断の像の全ての対称形式に関して自己共役な$g$は
$(gv)w-v(gw)\in K_E$を全ての$v,w$で満たす。
固有値$\lambda$の一つを取ると$\ker(g-\lambda)\cdot\operatorname{im}(g-\lambda)\subset K_E$
である。固有ベクトルの相手が高々一次元なので$\operatorname{rank}(g-\lambda)\le1$、
階数1なら像の生成元に三次元の相手があり矛盾する。従って$g$はスカラーである。
反対二位置で別の基底変換をして再び対称な像を得る場合、それらの相対変換は
全ての像形式に関し自己共役となる。もう一側にも同じ議論を適用する。
\end{proof}

\begin{corollary}[スキームとしての因子曲面の切断]
$\theta_x:w\mapsto xw\bmod K_E$の階数3以下の小行列式スキーム上で、
因子交換は$[x]\mapsto[\ker\theta_x]$という正則な対合になる。
左右のこのスキームは同型で、対合とも可換である。
\end{corollary}
\begin{proof}
階数2以下の点はないので三次小行列式の一つを可逆にすると、残る小行列式は
核直線が存在するためのスキーム論的方程式である。従ってその上の核直線のグラフが
順序付き因子スキームであり、補題\ref{lem:reciprocal-apolar-factors}を適用できる。
\end{proof}

\begin{theorem}[横断する二因子での階数の飛び]\label{thm:transverse-product-even-rank}
相異なる因子の対$x_0 x_1\in K_E$と対応する$f_0 f_1\in K_F$を取り、
四評価の階数を3とする。$U=U_{x_0}=V_{f_1}$、$V=U_{x_1}=V_{f_0}$が
相補的なら、反対階数は6または8である。特に両対称の階数7の標準AS候補では
\emph{全ての}因子対が一次元で交わり、その根基は可逆な二行二列関係を与える。
\end{theorem}
\begin{proof}
$U,V$は直交し各計量は非退化である。$e=\beta(x_0,f_1)\ne0$、
$d=\beta(x_1,f_0)\ne0$とする。全双線型写像は二つの積の和として
\[
 \beta(x,f)=(a(x)*_U b(f),\ m(x)*_V n(f))
\]
と書ける。ここで$a=\beta(-,f_1), b=\beta(x_0,-)$は$U$全体への写像であり、
それぞれ$kx_1,kf_0$を核とする。直交射影した$\beta$はこの二商を経由するので、
ある積$*_U$が一意に定まり、$e$は両側の単位である。同様に
$m=\beta(-,f_0),n=\beta(x_1,-)$と$V$の単位$d$を得る。
両対称式に$x_0$または$f_1$を代入すると、三形式
$\langle u*v,w\rangle_U$が全対称であることが分かる。
従って積は可換で計量は不変である。$V$でも同様で、結合性は仮定しない。

$D_U=U/ke$とし
\[
 T(p,q,r,s)=\langle p*q,r*s\rangle_U,\quad
 F(p,q;r,s)=T(p,r,q,s)-T(p,s,q,r)
\]
と書く。$F$は各対で交代し単位を入れると消えるため、
$D_U$の面積形式$\omega$を選べば$F=\kappa_U\omega(p,q)\omega(r,s)$である。
同じく$V/kd$で$\kappa_V$を取る。
$X=E^*/\langle x_0,x_1\rangle$、$Y=F^*/\langle f_0,f_1\rangle$からの誘導同型を
$A:X\to D_U,B:Y\to D_U,M:X\to D_V,N:Y\to D_V$と書く。
面積座標で両対称式は
\[
 \kappa_U\det A\det B+\kappa_V\det M\det N=0
\]
の一式となる。

反対行列から零行$x_0x_1$と零列$f_0f_1$を除く。
$x_0$の一次倍の三行と$f_1$の一次倍の三列は$U$の計量を与え、
$x_1$と$f_0$の三行三列は$V$の計量を与える。反対の組のブロックは零である。
この六行六列を消去した残差は$\operatorname{Sym}^2 X\times\operatorname{Sym}^2 Y$上にある。
一因子$U$の寄与は
\[
 \tfrac12\{T(p,r,q,s)+T(p,s,q,r)\}-T(p,q,r,s)
\]
の$(A,A;B,B)$による引き戻しである。実際、六行の消去は各積$p*q$を
単位を持つ三行から計量を介して差し引くことである。
同式は$-\tfrac12\{F(p,r;q,s)+F(p,s;q,r)\}$であり、
平方$(p,p;r,r)$で$-\kappa_U\omega(p,r)^2$となる。
二次元空間の対称二乗上の非退化対合$J$を
$J(p^2,r^2)=\omega(p,r)^2$で定めれば、残差はこの$J$の二つの引き戻しの線型結合である。

$R(g)=\operatorname{Sym}^2(g)/\det g$は$J$を保つ。二つの寄与の係数は
前の行列式の一式により逆符号なので、どちらも零なら残差は零である。
さもなくば残差の階数は$R(A^{-1}B)-R(M^{-1}N)$の階数である
（$X,Y,D_U,D_V$の基底を取って比較する）。
$R(g)$は面積形式による同定で無跡二行行列への共役作用になる。
二行非スカラー行列の中心化代数はその行列と恒等行列が張り、無跡部分は一次元である。
従って二つの$R$の差の階数は射影変換が等しければ0、異なれば2である。
六つのピボットを戻して結論を得る。最後は系\ref{cor:rank-seven-distinct-factors}による。
\end{proof}

\begin{lemma}[根基直線上の三次形式]\label{lem:product-radical-cubic}
同じ四軸で$U\cap V=kr$とし、根基の関係
$t(\bar x,\bar f)=\langle r,\beta(x,f)\rangle$は$X\times Y$上非退化とする。
$U$から定まる同型$J_U:X\to Y$を、$a=\beta(-,f_1)$と$b=\beta(x_0,-)$の
$U/k\beta(x_0,f_1)$への誘導同型により$\bar b^{-1}\bar a$と置く。
$V$からも同様に$J_V=\bar n^{-1}\bar m$を作る。
$J_U,J_V$が射影的に異なれば反対階数は8以上である。
射影的に等しければ階数は6または8で、6の場合には完全対称なテンソルに同値である。
\end{lemma}
\begin{proof}
$U$の元$u,b_1,b_2$に対して
$S_U(u,b_1,b_2)=\langle u,\beta(x,f)\rangle$を$a(x)=b_1,b(f)=b_2$で定める。
二つの持ち上げの曖昧さの像は$V=U^\perp$に入るため良定義である。
$x_0,f_1$を両対称式に代入すれば、第一と第三、第一と第二の引数交換を得る。
従って$S_U$は完全対称で、単位$e=\beta(x_0,f_1)$を入れると$U$の計量となる。
$S_U(r,-,-)$は$U/ke$上に降り、その$(\bar a,\bar b)$による引き戻しが$t$である。
$V$でも同様である。

前定理の六行六列のブロックは二つの退化計量で、各々階数2である。
四ピボットを消去する。$U$の根基行は$x_0 x_r$、$a(x_r)=r$で、
$\operatorname{Sym}^2Y$に対する残りの行は$S_U(r,\bar b(-),\bar b(-))$、
すなわち$t(J_U^{-1}-, -)$となる。根基列は$t(-,J_U-)$である。
$V$の根基行・列では$J_V$を使う。従って残る二行は$J_U,J_V$が射影的に
異なるとき独立で、二列も同時に独立である。根基行と根基列の交叉ブロックは零だから
全階数は$4+2+2$以上である。

射影的に等しければ$g:E^*\to F^*$を選び、$b g=c_U a,n g=c_V m$とできる
（$c_U,c_V\ne0$）。実際、商$X\to Y$を$c_UJ_U=c_VJ_V$と取り、
$x_0,x_1$を$c_U f_1,c_V f_0$へ送る。商の持ち上げは二つの係数を独立に直せ、
$b,n$の各単位方向のずれをちょうど消す。
この同定で$\beta(x,gy)$と$\beta(y,gx)$の差は$S_U,S_V$の対称性により
$U+V$と直交し、従って$kr$に入る。どちらかの文字が軸にあれば差は零である。
ゆえに交代部分は$r\eta$、$\eta$は二次元商$X$上の交代形式である。
$t_g(x,y)=t(x,gy)$もそこで対称かつ非退化である。

対称部分$\beta_s$だけのGramテンソル$w_s$は完全対称である。実際、
$\langle r,r\rangle=0$なので元のテンソルと同時転置との平均がそのGramテンソルであり、
両対称式を保つ。$\beta_s$の像は軸から$U+V$を含み、
$t_g\ne0$なので全六次元を張る。従って$w_s$の反対階数は6である。
残る交代寄与は$\operatorname{Sym}^2 X$同士にのみ入り、平方での値は
$2\eta(p,q)t_g(p,q)$である。四つの計量ピボットの後の根基行と根基列は
それぞれ$t_g$の非零倍一つに縮約される。従ってさらに二ピボットを取った後、
変化するのは$\ker(t_g:\operatorname{Sym}^2 X\to k)$の二次元上の交代対合だけである。
具体的にはこの核の各元を四つの軸ピボットと根基行の一つを使って持ち上げられ、
交代寄与は持ち上げで変わらない。$w_s$の残差は階数6から零である。
$t_g$を単位行列とする基底では、二つの無跡元$e_1^2-e_2^2,e_1e_2$の対合が
$2\eta(e_1,e_2)$になる。従って$\eta=0$で階数6、そうでなければ8である。
\end{proof}

\begin{theorem}[両対称AS窓の階数二分法]\label{thm:bisymmetric-six-or-eight}
標準14641型の両対称AS反転窓の反対階数は6または8であり、6の場合は
完全対称な窓に同値である。従ってこの類の低階数部分は対称完全交叉の分類に帰着し、
残る階数8の両核はともに四次完全交叉を定める。
\end{theorem}
\begin{proof}
楕円閉包定理で階数は8以下である。階数7なら系\ref{cor:rank-seven-distinct-factors}から
相異なる因子があり、交わりの次元は0または1である。定理\ref{thm:transverse-product-even-rank}と
補題\ref{lem:product-radical-cubic}でいずれも不可能である。
階数6以下なら$\mathbb P(K_E)$は次元3以上で、$\mathbb P^9$内の
次元6の階数2以下の対称二次式の多様体と交わる。
重平面なら定理\ref{thm:double-plane-ore}の$P=0$の場合で完全対称となる。
相異なる因子なら前二命題により階数は6以上であり、根基直線の場合の等号は既に示した。

横断する場合には$\kappa_U=0$なら$U$の可換積が結合的となる。
実際、$F=0$と計量の不変性・非退化性から二通りの三重積は一致する。
有限次元可換単位代数の指標$\chi:U\to k$を取ると、$\chi$を$V$上で零に延ばして
$\beta$と合成した関係は$\chi(a(x))\chi(b(f))$という非零純テンソルで、AS性に反する。
従って$\kappa_U\kappa_V\ne0$で、階数6なら前定理の計算から
$B^{-1}A$と$N^{-1}M$は射影的に一致する。前補題と同じ持ち上げで
$b g=c_U a,n g=c_V m$と取れる。両方の積は可換なので$\beta(x,gy)$は対称となり、
Gramテンソルは完全対称になる。最後の完全交叉性は定理\ref{thm:unique-linear-factor-partner}による。
\end{proof}
